\scanpage{22}
% Reused typescript, upright in the scan.
\begin{center}-2.24-\end{center}
of the terminology ``flag complex''.

A second direction in which we can make explicit this structure on $\Delta(G)$ is to show that the skeletons are pasted together in a sufficiently nice way so that various constructions which can be carried out in each skeleton, can also be carried out in $\Delta(G)$. The principal one we have in mind is drawing a great circle through two points. One approach to this is to complete the complex $\Delta(G)$ with respect to the metric $\rho$ mentioned above: then Tits has shown that the geodesic segments on this completion are precisely the segments of great circles of length $\leq\pi$, on the true spheres obtained by completing the skeletons $\Delta(T)$. Since we are avoiding this metric, we define these lines, and other auxiliary concepts, directly:

\uline{Definition 2.8}: A pair of points $\delta_1,\delta_2\in\Delta(T)$ is \uline{antipodal} if there is a 1-PS $\lambda$ of $G$ such that $\delta_1=\Delta(\lambda)$ and $\delta_2=\Delta(\lambda^{-1})$.

If $G$ is a torus, then every point of $\Delta(G)$ has a unique antipodal point, i.e.\ its antipodal on the sphere is the usual sense. More general, if $\delta_1$ is in the skeleton $\Delta(T)$ of $\Delta(G)$, for some maximal torus $T$, then the antipodal $\delta_2$ of $\delta_1$ as a point of $\Delta(T)$, is one of its antipodal points in $\Delta(G)$. In fact, these are the only points antipodal to $\delta_1$:

\uline{lemma 2.12}: If $\delta_1,\delta_2$ are two points in the skeleton $\Delta(T)$ of $\Delta(G)$, for some maximal torus $T$, then $\delta_1$ and $\delta_2$ are antipodal as points of $\Delta(G)$ if and only if they are antipodal as points of the sphere $\Delta(T)$.
\mspage{21}
